\section{Application to open problems using hybrid SAT-GA}\label{sec:app_hybrid_app}

This appendix details the specific experimental results and structural properties discovered by SolEvolve's Hybrid SAT-GA methodology when applied to open problems in coding theory, expanding upon the summary provided in Section~\ref{sec:hybrid}.

\subsection{Discovery of inequivalent $[35,10,12]$ codes}
A particularly significant contribution is the discovery of \textbf{10 mutually inequivalent} $[35,10,12]$ codes. While the Best Known Linear Code (BKLC) for parameters $[35,10]$ achieves minimum distance $d=12$, SolEvolve generated multiple codes matching this optimal distance but with distinct structural properties. We verified inequivalence using Magma's \texttt{IsEquivalent()} function.

Table~\ref{tab:inequiv35} summarizes three representative codes from our collection, each exhibiting a unique weight distribution.

\begin{table}[h!]
\centering
\caption{Inequivalent binary $[35,10,12]$ codes discovered by SolEvolve. All achieve the optimal BKLC distance but with distinct weight distributions. $A_i$ denotes the number of codewords of Hamming weight $i$ in the code.}
\label{tab:inequiv35}
\resizebox{\columnwidth}{!}{%
\begin{tabular}{@{}lccccccccc@{}}
\toprule
Code & $A_{12}$ & $A_{14}$ & $A_{16}$ & $A_{18}$ & $A_{20}$ & $A_{22}$ & $A_{24}$ & $A_{26}$ & $A_{28}$ \\ \midrule
BKLC & 180 & 0 & 549 & 0 & 276 & 0 & 18 & 0 & 0 \\
Sol-1 & 68 & 134 & 233 & 282 & 188 & 90 & 22 & 6 & 0 \\
Sol-5 & 98 & 161 & 245 & 255 & 186 & 63 & 14 & 1 & 0 \\
Sol-7 & 64 & 160 & 218 & 262 & 199 & 84 & 29 & 6 & 1 \\
\bottomrule
\end{tabular}%
}
\end{table}

\noindent The practical significance of discovering inequivalent codes extends beyond theoretical interest. Different codes with identical minimum distance may exhibit varying performance characteristics.

\subsection{$A_d$ optimization for $[43,10,16]$ codes}
Beyond finding inequivalent codes, SolEvolve demonstrates the capability for \textit{$A_d$ optimization}---iteratively reducing the number of minimum-weight codewords. Table~\ref{tab:ad_optim} shows the progression of discovered $[43,10,16]$ codes with decreasing $A_{16}$ values. Starting from the BKLC baseline of $A_{16} = 225$, the system iteratively discovers codes with fewer minimum-weight codewords, culminating in a code with $A_{16} = 80$---a 64\% reduction.

\begin{table}[h!]
\centering
\caption{$A_d$ optimization for $[43,10,16]$ codes. $A_i$ denotes the count of weight-$i$ codewords; minimizing $A_d$ improves undetected error probability. SolEvolve achieves a 64\% reduction from BKLC.}
\label{tab:ad_optim}
\resizebox{\columnwidth}{!}{%
\begin{tabular}{@{}lccccccccc@{}}
\toprule
Code & $A_{16}$ & $A_{18}$ & $A_{20}$ & $A_{22}$ & $A_{24}$ & $A_{26}$ & $A_{28}$ & $A_{30}$ & $A_{32}$ \\ \midrule
BKLC & 225 & --- & 472 & --- & 294 & --- & 32 & --- & --- \\
Sol-2 & 91 & 166 & 208 & 244 & 178 & 94 & 44 & 8 & 1 \\
Sol-3 & 86 & 166 & 210 & 247 & 178 & 90 & 44 & 8 & 1 \\
Sol-6 & \textbf{80} & 166 & 208 & 244 & 178 & 94 & 44 & 8 & 1 \\
\bottomrule
\end{tabular}%
}
\end{table}

\noindent All codes in Table~\ref{tab:ad_optim} were verified as inequivalent to the BKLC using Magma's \texttt{IsEquivalent()} function, confirming they represent genuinely new constructions rather than permutation equivalents.

\subsection{Matrix representation of the $[32,14,8]$ code}
The discovered $[32,14,8]$ code serves as a prime example of the system's capability. The generator matrix $G$ is presented below in standard form $G = [I_{14} \mid P]$:
\begin{center}
\resizebox{\linewidth}{!}{%
$G = \left[
\begin{array}{c|cccccccccccccccccc}
 & 1&0&1&1&0&0&1&1&0&1&1&1&0&0&0&1&0&1 \\
 & 0&1&0&1&0&0&1&0&1&1&1&0&1&0&0&1&1&0 \\
 & 1&1&1&1&0&1&0&1&0&0&1&0&1&0&1&0&1&1 \\
 & 1&1&1&1&0&1&0&0&1&0&1&0&1&0&0&1&0&0 \\
 & 0&1&1&0&0&1&1&0&1&0&1&0&1&1&1&0&1&1 \\
 & 1&0&0&0&0&0&0&0&1&0&1&0&1&0&1&0&1&1 \\
I_{14} & 1&1&1&1&0&0&1&1&0&0&0&0&1&1&0&0&0&1 \\
 & 0&1&1&1&1&0&1&0&1&0&1&0&1&1&0&1&0&1 \\
 & 1&0&0&1&1&1&1&0&1&0&1&0&1&1&1&1&1&1 \\
 & 0&1&0&1&1&0&0&1&0&0&1&1&0&1&1&1&0&1 \\
 & 0&1&1&0&0&0&0&0&0&0&0&0&1&1&0&1&1&1 \\
 & 1&0&1&1&0&1&1&1&1&0&1&0&0&1&1&0&0&1 \\
 & 0&1&0&0&0&0&1&1&0&0&1&0&0&0&1&1&1&0 \\
 & 1&1&0&0&1&0&0&1&0&0&0&0&0&1&1&0&1&0
\end{array}
\right]$%
}
\end{center}
The code's weight enumerator polynomial is:
\begin{align*}
W(z) =\; & 1 + 95z^8 + 54z^9 + 383z^{10} \\
&+ 310z^{11} + 1056z^{12} + 1044z^{13} \\
&+ 1960z^{14} + 2076z^{15} + 2223z^{16} \\
&+ 2368z^{17} + 1590z^{18} + 1576z^{19} \\
&+ 704z^{20} + 620z^{21} + 160z^{22} \\
&+ 132z^{23} + 17z^{24} + 10z^{25} \\
&+ 3z^{26} + 2z^{27}.
\end{align*}
